% !TeX root = ../../Book1.tex
\OptionalChapter{BV 函数的精细结构}{b1:ch:28}
\begin{chapterabstract}
本章把 BV 导数分成体积上的绝对连续部分、界面上的跳跃部分与剩余的 Cantor 部分。
通过近似上下极限、单侧值和跳跃法向，把等价类的点值恢复与测度分解联系起来。
最后介绍 SBV 空间，并在明确的值域与增长条件下推导 Mumford–Shah 型能量的极小元存在性。
精细迹、弥散结构及链律的深层输入单独标明，不以概览代替尚未展开的理论。
\end{chapterabstract}

\begin{learninggoals}
\begin{enumerate}
\item 用近似上下极限识别有限近似极限点，区分指定代表的近似连续与等价类自身的点态信息。
\item 确定跳跃的两侧值及法向，证明其唯一性与可测性，并用好层值解释跳跃集的可整流性。
\item 区分 \(D^au\)、\(D^ju\)、\(D^cu\)，辨认近似梯度、跳跃高度与奇异连续变化各自的作用。
\item 理解 SBV 为何排除 Cantor 部分，以及严格 BV 收敛为何仍不保证 SBV 闭合。
\item 在准确列出的深层输入下完成紧性、表面面积下半连续性和受限变分问题的存在性推导。
\end{enumerate}
\end{learninggoals}

上一章的余面积公式把一般函数还原为一族有限周长集合，
但这些层集的边界并不总是叠在同一张曲面上。
光滑函数的层面累积成体积中的梯度，阶跃函数的层面重叠成一道跳跃界面，
Cantor 函数则给出第三种可能：变化集中在零体积集合上，却没有任何非零跳跃。
导数测度的三部分分解正是对这三个现象的统一描述。

本章为选读，也是一卷实分析通向变分问题的收束。
近似极限的基础已在第十四章建立，这里增加上下极限与等价类代表的组织，
不重新改变“近似连续”对指定点值的含义。
读者应熟悉第六章的测度分解、第十九章的可整流性，以及上一章的 BV 逼近与余面积。
若暂时只想掌握可使用的结论，可以先看每个结构定理的准确陈述及一维例子，
再沿明确写出的证明概要追踪尚需学习的精细迹与切片工具。

Alberti 与 Mantegazza 的《A Note on the Theory of SBV Functions》%
\cite{AlbertiMantegazza1997SBV}提供了本章结构陈述与紧性路线的集中对照。
该文对基础精细结构另引专著，不能因为篇幅短就以为所有前置都已在几页内证明。
本书把能由已有工具推出的可整流性、测度分解后果、固定面积界下的闭合桥、
跳跃面积下半连续性和极小值论证写出；余下的深层输入明确标记。
进一步的系统阅读可选 Evans 与 Gariepy 的%
《Measure Theory and Fine Properties of Functions》\cite{Evans2015Measure}，
以及 Ambrosio、Fusco 与 Pallara 的%
《Functions of Bounded Variation and Free Discontinuity Problems》%
\cite{Ambrosio2000Functions}。

% !TeX root = ../../Book1.tex
\section{近似极限与近似极限代表}
\label{b1:sec:2801}

上一章把总变差解释为各层周长的累积。本章进一步问：这些周长在同一点附近怎样叠加？
有些点的函数值沿各个方向趋于同一个数，有些点则沿一个界面的两侧趋于不同的数。
若只查看任意选定代表的普通不连续点，零测修改就会掩盖这种区别。
因此先从邻近球中的体积比例定义点值，暂不使用导数测度。
本节的基础定义适用于一般实值可测函数；涉及有界变差函数的余维一结构时，再明确加入有限变差条件。

Alberti 与 Mantegazza 的《A Note on the Theory of SBV Functions》%
\cite[Remark 1.2]{AlbertiMantegazza1997SBV}把近似极限与导数分解放在一起介绍。
其文中把没有近似极限的集合记为 \(S_u\)，并在这个集合上给代表补零；
本章保留 \(S_u\) 表示近似不连续集，另以 \(J_u\) 专指有两个不同单侧值的跳跃点。
两者的关系是后面的结构定理，不能在定义时直接把它们当成同一个集合。
系统学习可继续参阅 Evans 与 Gariepy 的《Measure Theory and Fine Properties of Functions》%
\cite{Evans2015Measure}及 Ambrosio、Fusco 与 Pallara 的%
《Functions of Bounded Variation and Free Discontinuity Problems》%
\cite{Ambrosio2000Functions}。

% 实读AlbertiMantegazza1997SBV作者公开正式副本，刊印376--377页注1.2；
% 作者副本内页2--3。本文前节基础密度刻画独立证明，不冒称这些页含全证。
% 原文S_u与本书J_u的区别在正文说明；下节引述精细结构时另列证明义务。

\subsection{近似上、下极限}
\label{b1:subsec:280101}

固定开集 \(\Omega\subseteq\mathbb R^d\)，令 \(u:\Omega\to\mathbb R\) 可测。
只考察中心 \(x\in\Omega\) 且足够小的球。对可测集 \(A\subseteq\Omega\)，简记
\[
 \theta_r(A,x)=\frac{m_d(A\cap B(x,r))}{\omega_dr^d}.
\]
说一个性质在 \(x\) 附近以密度一成立，是说不满足它的点集之 \(\theta_r\) 趋于零。
这里的“密度一”不是“某个邻域中处处成立”，也不是“除掉一个固定的零测集后成立”。

\begin{definition}\label{b1:def:approximate-upper-lower-limits}
定义 \(u\) 在 \(x\) 的\term[jinsi-xia-jixian]{近似下极限}[approximate lower limit]与%
\term[jinsi-shang-jixian]{近似上极限}[approximate upper limit]为
\[
 u^\wedge(x)=
 \sup\{\,t\in\mathbb R:\lim_{r\downarrow0}\theta_r(\{u<t\},x)=0\,\},
\]
\[
 u^\vee(x)=
 \inf\{\,t\in\mathbb R:\lim_{r\downarrow0}\theta_r(\{u>t\},x)=0\,\}.
\]
允许取扩展实数值，约定 \(\sup\varnothing=-\infty\)、\(\inf\varnothing=+\infty\)。
\end{definition}
\symidx{\(u^\wedge\)：近似下极限}
\symidx{\(u^\vee\)：近似上极限}

下极限的容许阈值构成向下封闭的区间，上极限的容许阈值构成向上封闭的区间；
端点是否属于该区间不影响定义。不能把这里的体积密度换成小球中的本性上、下确界。
前者允许例外部分的比例逐渐趋零，后者只忽略严格零测的部分。

\begin{proposition}\label{b1:prop:approximate-limits-basic}
函数 \(u^\wedge,u^\vee\) 是扩展实值 Borel 函数，且
\[
 u^\wedge(x)\le u^\vee(x)\quad(x\in\Omega).
\]
它们只依赖 \(u\) 的几乎处处等价类。若 \(u\le v\) 几乎处处，则在每一点都有
\[
 u^\wedge\le v^\wedge,\quad u^\vee\le v^\vee.
\]
此外，\((-u)^\wedge=-u^\vee\)、\((-u)^\vee=-u^\wedge\)。
\end{proposition}
\begin{proof}
由\cref{b1:thm:lebesgue-borel-representative}，先把 \(u\) 换成几乎处处相等的 Borel 代表。
任意这样的替换都只在零测集内改变每个阈值集，
故不改变任何球中的体积，特别不改变任一点的两个近似极限。

对固定阈值 \(t\) 与半径 \(r\)，球体积积分作为中心的函数是 Borel 的。
对固定中心，它关于正半径连续，因为球壳的体积趋于零。
因此“\(\theta_r(\{u<t\},x)\to0\)”可以只用正有理半径写成可数个 Borel 条件。
例如对充分小的半径，使用
\(\inf_{n}\sup_{r\in\mathbb Q,\ 0<r<1/n}\theta_r\)；
球未必留在 \(\Omega\) 时可先在域外把阈值集补空，极限不受影响。
阈值的单调性又使定义中的上、下确界可只取有理 \(t\)。
所以两个近似极限都是 Borel 函数。

若 \(u^\wedge(x)>u^\vee(x)\)，可在二者之间取 \(a<b\)，使
\(\{u<b\}\) 与 \(\{u>a\}\) 在 \(x\) 都有密度零。
但这两个集合的并为 \(\Omega\)，不可能在内部一点有密度零，矛盾。
这个论证同样涵盖端点为无穷的情况。

若 \(u\le v\) 几乎处处，则
\(\{v<t\}\subseteq\{u<t\}\)、\(\{u>t\}\subseteq\{v>t\}\)，均按零测集忽略理解。
比较容许阈值区间即可得到两个单调关系。
最后，把 \(\{-u<t\}\) 写成 \(\{u>-t\}\)，便得到取负号的公式。
\end{proof}

\subsection{近似连续点}
\label{b1:subsec:280102}

\begin{definition}\label{b1:def:bv-finite-approximate-limit}
若存在有限实数 \(\ell\)，使对每个 \(\varepsilon>0\)，都有
\begin{equation}\label{b1:eq:approximate-limit-density}
 \lim_{r\downarrow0}
 \frac{m_d\bigl(\{y\in B(x,r):|u(y)-\ell|>\varepsilon\}\bigr)}
 {\omega_dr^d}=0,
\end{equation}
则按\cref{b1:def:approximate-limit}，\(u\) 在 \(x\) 有近似极限 \(\ell\)。
记该点集为 \(G_u\)，并在其上记 \(\widetilde u(x)=\ell\)。
\end{definition}

这里并不要求原代表已经满足 \(u(x)=\ell\)。
\chapref{b1:ch:14}的“近似连续”还要求这个等式成立，属于指定代表的性质；
本章的 \(G_u\) 只记录等价类有有限近似极限的点。
在 \(G_u\) 上重定义为 \(\widetilde u\) 后，该代表才在这些点近似连续。
后文为简便说等价类的近似连续点时，均指这项重定义后的性质，
不更改\chapref{b1:ch:14}对于指定函数值的定义。

\begin{proposition}\label{b1:prop:approximate-continuity-characterization}
对 \(x\in\Omega\) 及有限实数 \(\ell\)，下列条件等价：
\begin{equivalences}
 \item\label{b1:item:approximate-limit-density}
 \(u\) 在 \(x\) 的近似极限为 \(\ell\)。
 \item\label{b1:item:approximate-limits-equal}
 \(u^\wedge(x)=u^\vee(x)=\ell\)。
\end{equivalences}
近似极限若存在便唯一。若 \(u\in L^1_{\mathrm{loc}}(\Omega)\)，则几乎每点属于 \(G_u\)，
且 \(\widetilde u=u\) 几乎处处。
\end{proposition}
\begin{proof}
若密度条件成立，则对 \(a<\ell<b\)，集合 \(\{u<a\}\)、\(\{u>b\}\) 都有密度零。
故 \(u^\wedge\ge\ell\)、\(u^\vee\le\ell\)，结合前一命题得相等。
反过来，若两极限都等于 \(\ell\)，给定 \(\varepsilon>0\)，
可在 \(\ell-\varepsilon\) 与 \(\ell\) 之间取一个下极限的容许阈值，
在 \(\ell\) 与 \(\ell+\varepsilon\) 之间取一个上极限的容许阈值。
两项密度零条件的并包含 \(\{|u-\ell|>\varepsilon\}\)，所以该集合也有密度零。
唯一性随两极限的数值唯一性得到。

由\chapref{b1:ch:14}的 Lebesgue 微分定理，局部可积函数的几乎每点满足
\[
 \lim_{r\downarrow0}\frac1{\omega_dr^d}\int_{B(x,r)}|u(y)-u(x)|\,\mathrm dy=0.
\]
Chebyshev 不等式将左侧平均振荡控制转成%
\eqref{b1:eq:approximate-limit-density}，从而得到最后一项。
\end{proof}

Lebesgue 点的平均振荡趋零，比一般的密度意义近似极限强。
若 \(u\) 在 \(x\) 的某个邻域本性有界，二者确实等价：
密度意义收敛给出
\[
 \frac1{\omega_dr^d}\int_{B(x,r)}|u-\ell|
 \le\varepsilon+M\,\theta_r(\{|u-\ell|>\varepsilon\},x),
\]
其中 \(M\) 是邻域内 \(|u-\ell|\) 的本性上界。先令 \(r\downarrow0\)，再令
\(\varepsilon\downarrow0\) 即可。无界函数则可能在密度很小的部分上留下较大的积分。

\secref{b1:sec:1405}已构造可积尖峰的反例：非零部分在原点的密度趋零，
而球中归一化积分不趋零。因此这里不重复该构造，也不把近似极限存在自动当作一个
\(L^1\) 平均极限。有界变差函数在余维一尺度上还具有更强的性质，证明这点需要下一节的结构输入。

\subsection{不连续点集}
\label{b1:subsec:280103}

\begin{definition}\label{b1:def:bv-approximate-discontinuity-set}
定义 \(u\) 的\term[jinsi-bulianxu-ji]{近似不连续集}[approximate discontinuity set]为
\[
 S_u=\Omega\setminus G_u.
\]
\end{definition}
\symidx{\(S_u\)：没有有限近似极限的点集}
\symidx{\(\widetilde u\)：近似极限代表}

由前面的刻画，\(S_u\) 为 Borel 集。它包括 \(u^\wedge<u^\vee\) 的点，
也包括两极限都等于同一个无穷值的点；后者不能因上下极限相等而被算作有限值近似连续点。
在 \(S_u\) 上暂给 \(\widetilde u\) 补零，便得到一个全域 Borel 函数；
在尚未知道它与 \(u\) 几乎处处相等时，只把这项补值看作 Borel 规范化。
若 \(u\) 局部可积，前一命题进一步说明二者几乎处处相等；
此时才称 \(\widetilde u\) 为
\term[jinsi-jixian-daibiao]{近似极限代表}[approximate-limit representative]。
这个补零约定不声称零值具有特殊的边界意义；在跳跃点上，稍后还会定义两侧值的平均。

对任何 \(L^1_{\mathrm{loc}}\) 函数，目前只能由微分定理断言 \(m_d(S_u)=0\)。
对于有界变差函数，下一节将把它加强为一个余维一的结构结论。
加强的内容不是把 \(S_u\) 全部删掉，而是说明除了一个
\(\mathcal H^{d-1}\)-零集，近似不连续表现为沿界面的两侧跳跃。

例如 \(u=\mathbf1_{\{x_d>0\}}\) 在平面外近似连续，在平面上有
\(u^\wedge=0\)、\(u^\vee=1\)，因此 \(S_u=\{x_d=0\}\)。
它是局部有界变差函数，却在一个具有正余维一测度的集合上没有单一近似极限。
反之，\(\mathbf1_{\mathbb Q^d}\) 的任意点均有近似极限零，
所以 \(S_u=\varnothing\)，尽管这个指定代表在通常意义下处处不连续。
两个例子分别说明，跳跃不能通过给界面点重赋单一值消除，
而零测修改造成的普通不连续根本不应计入 \(S_u\)。

无穷近似值即使对有界变差函数也可能发生。设 \(d\ge2\)、\(0<\alpha<d-1\)，
在 \(B(0,1)\) 上令 \(u(x)=|x|^{-\alpha}\)，原点任意补值。
球壳积分给出 \(u,\,\alpha|x|^{-\alpha-1}\in L^1\)。
在删去半径 \(\varepsilon\) 小球的区域上分部积分，内边界项的绝对值至多为
\[
 C\varepsilon^{d-1-\alpha}\|\varphi\|_\infty\longrightarrow0.
\]
因此经典梯度在原点外的表达式也是全域弱梯度，\(u\in W^{1,1}\subset BV\)。
但原点的两个近似极限均为 \(+\infty\)，故 \(0\in S_u\)，而这不是两个有限值之间的跳跃。
单点的 \(\mathcal H^{d-1}\) 测度为零，所以这个现象与下一节的结构定理相容。

\begin{proposition}[近似振荡与层集边界]\label{b1:prop:approximate-oscillation-level-boundary}
设 \(T\subset\mathbb R\) 为任意可数稠密集。则
\[
 \{x:u^\wedge(x)<u^\vee(x)\}
 \subseteq\bigcup_{t\in T}\partial^e\{u>t\}.
\]
更具体地，若 \(u^\wedge(x)<t<u^\vee(x)\)，则 \(x\) 属于该层集的测度论边界；
若 \(t<u^\wedge(x)\)，层集在 \(x\) 的密度为 \(1\)；
若 \(t>u^\vee(x)\)，密度为 \(0\)。
\end{proposition}
\begin{proof}
若 \(t<u^\wedge(x)\)，在 \(t\) 与下极限之间选容许阈值 \(a\)，
则 \(\{u\le t\}\subseteq\{u<a\}\) 的密度为零，故 \(\{u>t\}\) 密度为一。
上极限情形直接由阈值单调性得到。

若 \(u^\wedge(x)<t<u^\vee(x)\)，而 \(\{u>t\}\) 密度为零，
则 \(t\) 是上极限的容许阈值，与 \(t<u^\vee(x)\) 矛盾。
若它密度为一，则 \(\{u<t\}\) 密度为零，推出 \(u^\wedge(x)\ge t\)，同样矛盾。
因而该点既不在测度论内部，也不在外部，属于 \(\partial^e\{u>t\}\)。
最后，在每段非空的扩展实数开区间 \((u^\wedge(x),u^\vee(x))\) 中取一个 \(T\) 的点，
就得到可数并包含。
\end{proof}

这条联系允许先选出一组周长有限的稠密好层值，再利用上一章的边界结构。
它还没有处理两极限同为无穷的点，也没有保证同一点上的不同层面法向一致；
后两项正是不能由这个简单包含替代的精细性质。

后续所有界面与代表仍由几乎处处类决定。特别是，
“导数由函数的精确代表给出”必须同时说明用了哪一种代表、
在哪一个测度的几乎处处意义下成立；Lebesgue 几乎处处相等的代表，
若被直接放进奇异导数测度的积分，结果未必相同。

% !TeX root = ../../Book1.tex
\section{跳跃集}
\label{b1:sec:2802}

近似不连续只说明一个球中不能用单一数值描述函数。
跳跃则有更确定的形状：把球沿一个平面分成两半后，两侧分别接近两个不同的有限值。
本节先把这个局部图像写成定义，确定两侧值和方向；
再用上一章的层集与周长结构，说明有界变差函数的跳跃点为何由可数个曲面承载。

以下均取实值函数。定义和最初几个命题只要求 \(u\in L^1_{\mathrm{loc}}(\Omega)\)；
涉及可整流性与导数测度时，另加 \(u\in BV(\Omega)\)。
文献有时允许任意选择法向并同时交换两侧值，本书则把较大值记作 \(u^+\)，固定取向。
这里的上标 \(u^\pm\) 专用于跳跃迹；一般函数的正部、负部已统一记为
\(u_+\)、\(u_-\)，二者不得混用。
这一选择不改变最终的向量测度。

% 实际阅读：AlbertiMantegazza1997SBV，作者公开正式副本，注1.2，
% 作者版内页2--3（刊印376--377页）。其中给出广义跳跃集上的两侧L1迹与导数公式，
% 并将精细结构证明指向Evans--Gariepy 1992第5.9节；未声称已实读该书证明。
% 下文唯一性、上下值关系、Borel性及Ju可整流性在本书工具内独立证明；
% 非有限近似值的余维一零测性、BV曲面两侧迹与界面导数公式明确作为深层输入。

\subsection{跳跃点}
\label{b1:subsec:280201}

对 \(x\in\Omega\)、单位向量 \(\nu\) 及足够小的 \(r>0\)，记
\[
 B_r^\pm(x,\nu)
 =\{y\in B(x,r):\ \pm(y-x)\cdot\nu>0\}.
\]
两半球的体积均为 \(\omega_dr^d/2\)；中间平面为体积零集，不影响平均。

\begin{definition}\label{b1:def:approximate-jump-point}
给定 \(u\in L^1_{\mathrm{loc}}(\Omega)\) 与 \(x\in\Omega\)。
若存在有限实数 \(a<b\) 和单位向量 \(\nu\)，使
\begin{equation}\label{b1:eq:approximate-jump-halves}
 \begin{aligned}
 \lim_{r\downarrow0}\frac2{\omega_dr^d}
       \int_{B_r^+(x,\nu)}|u(y)-b|\,\mathrm dy&=0,\\
 \lim_{r\downarrow0}\frac2{\omega_dr^d}
       \int_{B_r^-(x,\nu)}|u(y)-a|\,\mathrm dy&=0.
 \end{aligned}
\end{equation}
则称 \(x\) 为 \(u\) 的\term[jinsi-tiaoyue-dian]{近似跳跃点}[approximate jump point]。
全体这样的点组成\term[tiaoyue-ji]{跳跃集}[jump set]，记为 \(J_u\)。
\end{definition}
\symidx{\(J_u\)：函数的近似跳跃点集}

此处使用的是半球中的 \(L^1\) 平均，不只是某个稀疏点列上的极限，
也不只是半球内依测度趋近一个数。尤其对无界函数，
上一节已经说明后两种较弱信息不能自动给出平均绝对误差趋零。

把半球缩放到单位球更便于比较不同的跳跃描述。设
\[
 w_{a,b,\nu}(z)=
 \begin{cases}
  b,&z\cdot\nu>0,\\
  a,&z\cdot\nu<0.
 \end{cases}
\]
平面上任意补值。条件\eqref{b1:eq:approximate-jump-halves}恰好等价于
\begin{equation}\label{b1:eq:jump-step-blowup}
 u(x+r\,\cdot)\longrightarrow w_{a,b,\nu}
 \quad\text{于 }L^1\bigl(B(0,1)\bigr).
\end{equation}
这个缩放极限保留了跳跃的高度和方向，却不读取原代表在点 \(x\) 的指定值。

\begin{proposition}\label{b1:prop:approximate-jump-uniqueness}
在 \(a<b\) 的排序约定下，近似跳跃点的三元组 \((a,b,\nu)\) 唯一。
跳跃集及其三元组只依赖函数的几乎处处等价类。
\end{proposition}
\begin{proof}
若两个三元组满足定义，则由 \(L^1\) 极限的唯一性，所得两个阶跃函数几乎处处相等。
每个阶跃函数都恰好取两个不同的本性值，而且每个值占半个单位球。
因此它们的无序值集相同；排序以后，两侧的数值分别相同。

较大值所占的半球于是也几乎处处相同。若两个单位法向不同，
对应正半球的对称差便含有一个非空开锥与单位球的交，具有正体积，矛盾。
故法向也相同。若不排序，唯一的另一种写法是同时交换两个值并把法向反号。
最后，零测修改不改变任何半球积分，故不改变定义。
\end{proof}

\subsection{单侧近似极限}
\label{b1:subsec:280202}

现在可以在 \(J_u\) 上无歧义地记
\[
 u^-(x)\coloneqq a,\quad u^+(x)\coloneqq b,\quad \nu_u(x)\coloneqq\nu.
\]
称 \(u^\pm(x)\) 为相应两侧的\term[dance-jinsi-jixian]{单侧近似极限}[one-sided approximate limits]。
这里的正负号指向由 \(\nu_u\) 确定的半球；依本书约定，正侧同时是数值较大的一侧。
\symidx{\(u^+\)：跳跃点的较大单侧近似极限}
\symidx{\(u^-\)：跳跃点的较小单侧近似极限}

\begin{proposition}\label{b1:prop:jump-approximate-limits}
若 \(x\in J_u\)，则
\begin{equation}\label{b1:eq:jump-upper-lower-values}
 u^\wedge(x)=u^-(x)<u^+(x)=u^\vee(x).
\end{equation}
特别地，\(J_u\subseteq S_u\)。对每个 \(u^-(x)<t<u^+(x)\)，还有
\begin{equation}\label{b1:eq:jump-level-half-density}
 \lim_{r\downarrow0}
 \frac{m_d\bigl(\{u>t\}\cap B(x,r)\bigr)}{\omega_dr^d}
 =\frac12.
\end{equation}
\end{proposition}
\begin{proof}
记 \(a=u^-(x)\)、\(b=u^+(x)\)。在正半球中，\(\{u\le t\}\) 上的
\(|u-b|\) 至少为 \(b-t>0\)，故 Chebyshev 不等式说明这个例外集的相对体积趋零。
在负半球中同理，\(\{u>t\}\) 的相对体积趋零。
两项相加就得到\eqref{b1:eq:jump-level-half-density}。

若 \(s<a\)，两半球中的误差估计都给出 \(\{u<s\}\) 的密度为零；
若 \(s>a\)，在 \(a\) 与 \(\min(s,b)\) 之间取一个阈值，
负半球便给出 \(\{u<s\}\) 的正密度下界。
所以近似下极限恰为 \(a\)。对上极限作相同的论证，得到 \(b\)。
由于两者不同，有限近似极限不可能存在，故 \(x\in S_u\)。
\end{proof}

在跳跃点上还可取中点值
\[
 u^*(x)=\frac{u^+(x)+u^-(x)}2.
\]
由两个半球的平均收敛立即得到
\[
 \lim_{r\downarrow0}\frac1{\omega_dr^d}\int_{B(x,r)}u(y)\,\mathrm dy=u^*(x).
\]
但这不使跳跃点成为 Lebesgue 点：事实上
\[
 \lim_{r\downarrow0}\frac1{\omega_dr^d}\int_{B(x,r)}|u(y)-u^*(x)|\,\mathrm dy
 =\frac{u^+(x)-u^-(x)}2>0.
\]
因此球平均有极限、有限近似极限存在、平均绝对振荡趋零，是需要区别的三种陈述。
因此这里的 \(u^*\) 与\chapref{b1:ch:14}由球平均定义的精确代表在 \(J_u\) 上相同；
本书始终把 \(u^*\) 称为球平均精确代表，而把 \(\widetilde u\) 称为近似极限代表，
不另用这两个记号表示彼此冲突的全域点值约定。

\subsection{跳跃法向量}
\label{b1:subsec:280203}

向量 \(\nu_u\) 称为\term[tiaoyue-faxiangliang]{跳跃法向量}[jump normal]。
它从低值侧指向高值侧。这个取向首先由函数的两侧行为定义，
尚不需要预先知道 \(J_u\) 是一张光滑曲面。
\symidx{\(\nu_u\)：从低值侧指向高值侧的单位跳跃法向量}

取盒子 \(\Omega=(-1,1)^d\)，并令
\[
 u(x)=
 \begin{cases}
  a,&x_d<0,\\
  b,&x_d>0,
 \end{cases}
 \quad a<b.
\]
界面上任意补值。此时
\[
 J_u=\Omega\cap\{x_d=0\},\quad
 u^-=a,\quad u^+=b,\quad \nu_u=e_d.
\]
逐坐标应用 Fubini 定理和一维分部积分，可得
\[
 Du=(b-a)e_d\,\mathcal H^{d-1}\llcorner
       \bigl(\Omega\cap\{x_d=0\}\bigr).
\]
具体地，前 \(d-1\) 个方向的测试积分为零；最后一个方向则满足
\[
 -\int_\Omega u\,\partial_d\varphi\,\mathrm dx
 =(b-a)\int_{(-1,1)^{d-1}}\varphi(x',0)\,\mathrm dx'.
\]
这也直接检验了导数中法向因子的正负号。

对光滑集合的示性函数，数值 \(1\) 位于集合内部，数值 \(0\) 位于外部。
所以在其边界上，\(\nu_{\mathbf1_E}=-\nu_E\)，其中 \(\nu_E\) 是集合外法向。
跳跃公式因此写成
\[
 D\mathbf1_E=-\nu_E\,\mathcal H^{d-1}\llcorner\partial E,
\]
与\cref{b1:prop:smooth-domain-perimeter}一致。不能将集合外法向直接代入
排序后的函数跳跃公式而漏掉这个负号。

为了稍后把两侧值和法向放进测度积分，还需确认它们的可测性。
\begin{proposition}\label{b1:prop:jump-data-borel}
对 \(u\in L^1_{\mathrm{loc}}(\Omega)\)，集合 \(J_u\) 是 Borel 集，
而 \(u^\pm\) 与 \(\nu_u\) 是其上的 Borel 函数。
\end{proposition}
\begin{proof}
由\cref{b1:thm:lebesgue-borel-representative}，先选取 \(u\) 的 Borel 代表；
这不改变任何跳跃数据。
由\cref{b1:prop:approximate-limits-basic}，集合
\[
 A=\{x:-\infty<u^\wedge(x)<u^\vee(x)<+\infty\}
\]
是 Borel 的。在 \(A\) 上置 \(a=u^\wedge(x)\)、\(b=u^\vee(x)\)，并记
\[
 F_r(x,v)=\frac1{\omega_d}\int_{B(0,1)}
       |u(x+rz)-w_{a,b,v}(z)|\,\mathrm dz.
\]
先在距边界有正距离的可数个内部集合上考察足够小的 \(r\)。
对固定 \(r,v\)，这个函数关于 \(x\) 是 Borel 的；关于正半径 \(r\) 连续。
后一事实先对连续函数成立，再在共同紧集上用 \(L^1\) 逼近即可得到。
它关于 \(v\) 也连续，因为两半球的对称差体积随方向连续变化。

取单位球面的一个可数稠密子集 \(D\)。条件 \(x\in J_u\) 等价于：
对每个正整数 \(k\)，存在 \(v\in D\) 和正整数 \(m\)，使所有充分留在定义域内、
满足 \(0<r<1/m\) 的正有理半径都有 \(F_r(x,v)<1/k\)。
这是一组可数的 Borel 条件。
正向蕴含由阶跃极限及方向的连续性立即得到。
反过来，选取这些方向 \(v_k\)；在同时足够小的半径上用三角不等式，得到
\[
 \frac1{\omega_d}\int_{B(0,1)}
 |w_{a,b,v_k}-w_{a,b,v_\ell}|
 \le \frac1k+\frac1\ell.
\]
从单位球面到这些阶跃函数的映射连续且单射，球面又紧，
故 \(v_k\) 收敛到唯一方向 \(v\)。再次用三角不等式，并利用半径连续性，
便得 \(F_r(x,v)\to0\)。因此上述可数条件确实刻画 \(J_u\)。

每次按固定枚举选取第一个合条件的 \((v,m)\)，便使 \(v_k(x)\) 成为 Borel 函数；
其极限就是 \(\nu_u(x)\)。最后由\eqref{b1:eq:jump-upper-lower-values}，
两侧值也是 Borel 函数。
\end{proof}

\subsection{\texorpdfstring{\(J_u\)}{Ju} 的可整流性}
\label{b1:subsec:280204}

跳跃方向并不预设曲面，但有界变差条件将这些局部方向约束在可数个可整流集合上。
证明的关键是只选一组稠密的好层值，而不要求每一个层值都有有限周长。

\begin{proposition}\label{b1:prop:bv-jump-rectifiable}
若 \(u\in BV(\Omega)\)，则 \(J_u\) 可数 \((d-1)\)-可整流。
\end{proposition}
\begin{proof}
由有界变差函数的余面积公式，几乎每个 \(t\in\mathbb R\) 都使
\(E_t=\{u>t\}\) 在 \(\Omega\) 内有有限周长。
这个好层值集合中可以选取一个可数稠密集 \(D\)：在每个有理端点开区间中取一个好值即可。
不能未经选择便声称所有有理层值都好。

若 \(x\in J_u\)，在 \((u^-(x),u^+(x))\) 内取 \(t\in D\)。
由\eqref{b1:eq:jump-level-half-density}，\(E_t\) 在 \(x\) 的密度为 \(1/2\)，
所以 \(x\in\partial^eE_t\)。于是
\[
 J_u\subseteq\bigcup_{t\in D}(\partial^eE_t\cap\Omega).
\]
De Giorgi 结构\cref{b1:thm:de-giorgi-structure}说明，
每个 \(\partial^eE_t\cap\Omega\) 与可数可整流的
\(\partial^*E_t\cap\Omega\) 只差一个 \(\mathcal H^{d-1}\)-零集。
可数合并后就得到所需覆盖。这里没有声称每个跳跃点都属于某个指定层集的约化边界；
用到的是两类边界之间已经证明的模零关系。
当 \(d=1\) 时，各好层集的内部边界局部只有有限个点，
同一论证给出 \(J_u\) 至多可数。
\end{proof}

可整流性尚未回答两个问题：近似不连续能否在跳跃集之外占据正的余维一测度，
以及导数在跳跃面上究竟有多大。完整答案如下。
\begin{theorem}[有界变差函数的跳跃结构（深层结构定理）]\label{b1:thm:bv-jump-structure}
若 \(u\in BV(\Omega)\)，则
\begin{equation}\label{b1:eq:bv-discontinuity-jump-null}
 \mathcal H^{d-1}(S_u\setminus J_u)=0.
\end{equation}
而且其向量导数测度在 \(J_u\) 上满足
\begin{equation}\label{b1:eq:bv-jump-derivative}
 Du\llcorner J_u
 =(u^+-u^-)\nu_u\,\mathcal H^{d-1}\llcorner J_u.
\end{equation}
因而
\[
 |Du|\llcorner J_u
 =(u^+-u^-)\,\mathcal H^{d-1}\llcorner J_u,
 \quad
 \int_{J_u}(u^+-u^-)\,\mathrm d\mathcal H^{d-1}
 \le |Du|(\Omega).
\]
\end{theorem}
\begin{proofsketch}
上述可整流性已完整证明。余下的两项是有界变差函数精细结构的深层部分；
Alberti 与 Mantegazza 在《A Note on the Theory of SBV Functions》的注1.2中准确列出相应的两侧 \(L^1\) 迹及导数公式%
\cite[Remark 1.2, pp. 376--377]{AlbertiMantegazza1997SBV}。
该注把精细结构的证明指向 Evans 与 Gariepy 的相关论述，
并不在这几页中完整证明曲面迹定理。这里说明尚需的桥梁。

第一项输入是有界变差函数的非有限近似值集合
\[
 N_\infty
 =\{x:u^\wedge(x)\notin\mathbb R
          \text{ 或 }u^\vee(x)\notin\mathbb R\}
\]
具有零 \(\mathcal H^{d-1}\) 测度。
在 \(S_u\setminus N_\infty\) 上，两近似极限有限且不相等。
在它们之间选稠密好层值，由\cref{b1:prop:approximate-oscillation-level-boundary}，
便把这些点覆盖在好层集的测度论边界上。
由上一章，这些边界模零由可数个可整流曲面承载。

第二项输入是有界变差函数在可整流曲面上的两侧迹定理：
对曲面面积几乎每点，都存在有限的两个半球 \(L^1\) 迹；
在该曲面上的导数限制等于两侧迹之差乘有向单位法向与面积测度。
两迹相等时，合并半球平均便得到有限近似极限；
故在所覆盖的 \(S_u\) 上，两迹几乎处处不同，按大小排序即得到 \(J_u\)。
这证明\eqref{b1:eq:bv-discontinuity-jump-null}；
曲面上的导数限制公式再给出\eqref{b1:eq:bv-jump-derivative}。
最后，单位法向的长度为一，故向量测度的变差就是所写的非负密度积分。

本节不展开非有限近似值的余维一估计，以及有界变差函数在曲面两侧的迹及其分部积分公式。
它们需要进一步的有界变差函数切片和精细代表理论，不能把第23章的 Sobolev 迹定理
直接套在一般有界变差函数上来替代。
\end{proofsketch}

上一节的无穷近似值例子因此没有被漏掉：那些点仍属于 \(S_u\)，
却不符合两个有限值的跳跃定义；结构定理把它们连同其他非跳跃例外一起控制在
\(\mathcal H^{d-1}\)-零集中。
另一方面，公式控制的是跳跃高度与面积的乘积，
并不保证跳跃集本身有有限面积。
它只立即保证
\[
 \mathcal H^{d-1}
 \bigl(\{x\in J_u:u^+(x)-u^-(x)\ge 1/k\}\bigr)
 \le k\,|Du|(\Omega).
\]
下一节将把这部分导数与体积上的密度、以及剩余的 Cantor 部分分开，
从而得到三种互不混淆的变差来源。

% !TeX root = ../../Book1.tex
\section{导数测度的分解}
\label{b1:sec:2803}

分布导数把一个函数的全部一阶变化保存在同一向量测度中。
但这些变化可以分布在体积中，也可以集中在跳跃面上，还可能既不占体积、又不由跳跃面承载。
本节区分这三种情形。前两种分别连接 Sobolev 梯度与有限周长集合，
第三种则解释连续函数为何仍可能没有可积的弱导数。

以下 \(\Omega\subseteq\mathbb R^d\) 为开集，\(u\in BV(\Omega)\) 为实值。
跳跃值与法向沿用\secref{b1:sec:2802}的约定：
\(u^+>u^-\)，\(\nu_u\) 从低值侧指向高值侧。
深层结构的准确陈述可参看 Alberti、Mantegazza 的
《A Note on the Theory of SBV Functions》\cite[Remark 1.2]{AlbertiMantegazza1997SBV}；
系统学习可继续阅读 Ambrosio、Fusco、Pallara 的
《Functions of Bounded Variation and Free Discontinuity Problems》%
\cite[Chapter 3]{Ambrosio2000Functions}。
本节把测度分解的直接后果完整推出，而将近似微分的识别和链律中的深层步骤明确列为证明概要。
% 实读作者公开正式论文 Alberti--Mantegazza 1997：
% 作者版第2--3页（刊印376--377）Remark 1.2；
% 作者版第4--5页（刊印378--379）Remark 2.1；
% 其§2.3--2.5另供下一节紧性使用。AFP只核实出版与章节导航，无全文阅读声明。
% 原文Su表示无近似极限集；本书保留Su与真正两侧跳跃集Ju的区别。

\subsection{绝对连续部分}
\label{b1:subsec:280301}

先不使用函数的几何性质。对 \(Du\) 的每个分量应用 Lebesgue 分解，
得到唯一的表示
\begin{equation}\label{b1:eq:bv-lebesgue-decomposition}
 Du=g\,m_d+D^su,\quad g\in L^1(\Omega;\mathbb R^d),\quad
 |D^su|\perp m_d.
\end{equation}
记 \(D^au=g\,m_d\)，称它为导数测度的绝对连续部分。
具体说，各分量的奇异部分分别集中在某个体积零的 Borel 集上；
取这有限多个集合的并 \(N\)，便有
\[
 D^au=Du\mathbin{\llcorner}(\Omega\setminus N),
 \quad
 D^su=Du\mathbin{\llcorner}N.
\]
这里限制记号只表示在相应 Borel 集上取测度。
向量极分解给出 \(|D^au|=|g|m_d\)；
不同分量的分解由标量唯一性保证相容。因而这一步只需要第6章的测度论，
尚未证明 \(g\) 与逐点的微分有什么联系。
\symidx{\(D^au\)：\(Du\) 的绝对连续部分}
\symidx{\(D^su\)：\(Du\) 的奇异部分}

\begin{theorem}[有界变差函数的近似梯度]
\label{b1:thm:bv-approximate-gradient}
函数 \(u\) 在 \(m_d\)-几乎处处近似可微。
若在近似极限存在处用该极限选取中心值，则其近似微分可写成
\[
 \operatorname{ap}Du(x)\,h=\nabla u(x)\cdot h,
 \quad \nabla u(x)=g(x)
 \quad m_d\text{-几乎处处},
\]
其中 \(g\) 是\eqref{b1:eq:bv-lebesgue-decomposition}中的密度。
特别地，\(D^au=\nabla u\,m_d\)，且 \(\nabla u\in L^1(\Omega;\mathbb R^d)\)。
\end{theorem}
\begin{proofsketch}
这里不仅需要 Radon--Nikodym 定理，还要证明函数本身在几乎每点有一阶线性近似。
一种路线先用变差的局部平均控制函数振荡，把函数限制在可数个 Lipschitz 片上；
在这些片的密度点应用 Rademacher 定理，再用分布测试识别所得梯度与 \(g\)。
把振荡估计转成覆盖几乎所有点的 Lipschitz 分片，以及最后的密度识别，
是本节未展开的两步，不能由测度有密度直接代替。
上述结论在所引短文 Remark 1.2 和 equation (1.2) 之后明确陈述，
该处将完整论证指向 Evans--Gariepy 原版的 Section 6.1。
\end{proofsketch}

从此对一般有界变差函数写 \(\nabla u\) 时，指这个近似梯度的几乎处处等价类。
它不是 \(Du\) 的另一种写法，也不要求一个任意指定的代表具有经典偏导数。
若 \(u\in W^{1,1}(\Omega)\)，分布导数已经完全由弱梯度表示，
上述唯一性便说明近似梯度与弱梯度几乎处处一致。

\begin{proposition}\label{b1:prop:bv-sobolev-singular-criterion}
对于 \(u\in BV(\Omega)\)，有
\[
 u\in W^{1,1}(\Omega)\quad\Longleftrightarrow\quad D^su=0.
\]
\end{proposition}
\begin{proof}
正向由弱导数的定义和 Lebesgue 分解唯一性得到。
反向若 \(D^su=0\)，则每个分布偏导数都由 \(g_i\in L^1(\Omega)\) 表示。
再加上有界变差函数定义中的 \(u\in L^1(\Omega)\)，恰好满足 \(W^{1,1}\) 的定义。
\end{proof}

\subsection{跳跃部分}
\label{b1:subsec:280302}

\begin{definition}\label{b1:def:bv-jump-derivative}
将导数测度限制到跳跃集 \(J_u\)，所得向量测度记为
\[
 D^ju=Du\mathbin{\llcorner}J_u,
\]
称为导数测度的\term[tiaoyue-bufen]{跳跃部分}[jump part]。
\end{definition}
\symidx{\(D^ju\)：\(Du\) 的跳跃部分}

由\cref{b1:thm:bv-jump-structure}，它具有明确的表面密度：
\begin{equation}\label{b1:eq:bv-jump-density}
 D^ju=(u^+-u^-)\nu_u\,
       \mathcal H^{d-1}\mathbin{\llcorner}J_u.
\end{equation}
跳跃的高度、朝向与承载的面积在这个公式中各占一个位置。
如果改用不排序的两侧记号，同时交换 \(u^+\)、\(u^-\) 并把法向反号，
右侧向量不变。因此测度不依赖人为的定向选择。
由于 \(|\nu_u|=1\)，对每个 Borel 集 \(A\subseteq\Omega\)，
\[
 |D^ju|(A)=\int_{A\cap J_u}(u^+-u^-)\,\mathrm d\mathcal H^{d-1}.
\]
这给出跳跃高度的可积性，但不保证 \(\mathcal H^{d-1}(J_u)<\infty\)：
很多很小的跳跃，可以具有有限总变差和无限总面积。

例如在 \((0,1)\) 内令
\[
 v(x)=\sum_{n=1}^{\infty}2^{-n}\mathbf1_{(2^{-n},1)}(x).
\]
该级数一致绝对收敛，每项的分布导数为 \(2^{-n}\delta_{2^{-n}}\)。
对测试函数逐项积分合法，故
\[
 Dv=\sum_{n=1}^{\infty}2^{-n}\delta_{2^{-n}},
 \quad |Dv|\bigl((0,1)\bigr)=1.
\]
在每个 \(2^{-n}\) 处跳跃高度为 \(2^{-n}\)，其他内部点在一个邻域内函数恒定，
所以 \(J_v=\{\,2^{-n}:n\ge1\,\}\)。
一维的 \(\mathcal H^0\) 是计数测度，因而 \(\mathcal H^0(J_v)=\infty\)。

\subsection{Cantor 部分}
\label{b1:subsec:280303}

跳跃集可数 \((d-1)\)-可整流，因而是体积零集。
所以 \(D^ju\) 是 \(D^su\) 在 \(J_u\) 上的限制，余下的部分仍为奇异测度。

\begin{definition}\label{b1:def:bv-cantor-derivative}
规定
\[
 D^cu=D^su-D^ju
     =D^su\mathbin{\llcorner}(\Omega\setminus J_u),
\]
称 \(D^cu\) 为导数测度的\term[cantor-bufen]{Cantor 部分}[Cantor part]。
\end{definition}
\symidx{\(D^cu\)：\(Du\) 的 Cantor 部分}

\begin{theorem}[有界变差函数的导数测度三分解]
\label{b1:thm:bv-derivative-decomposition}
对于任意实值 \(u\in BV(\Omega)\)，有
\begin{equation}\label{b1:eq:bv-three-part-decomposition}
 Du=\nabla u\,m_d
    +(u^+-u^-)\nu_u\,
      \mathcal H^{d-1}\mathbin{\llcorner}J_u
    +D^cu.
\end{equation}
这三项两两互相奇异。更进一步，若 Borel 集 \(A\subseteq\Omega\)
关于 \(\mathcal H^{d-1}\) 为 \(\sigma\)-有限，则
\[
 |D^cu|(A)=0.
\]
有限近似极限 \(\widetilde u\) 在 \(|D^cu|\)-几乎处处存在。
\end{theorem}
\begin{proofsketch}
分解等式由前面两项的识别和 \(D^cu\) 的定义得到。
若 \(N\) 承载 \(D^su\) 且 \(m_d(N)=0\)，三项分别集中于
\(\Omega\setminus(N\cup J_u)\)、\(J_u\)、\(N\setminus J_u\)，故两两互相奇异。
这里尚不能只由 \(D^cu(J_u)=0\) 推出它不充任何有限面积集。

后一性质是有界变差函数的精细结构所给出的额外结论。
证明须把任意有限 \(\mathcal H^{d-1}\) 集与函数在不同方向上的一维截面比较，
识别其中能承载导数的部分为两侧迹的跳跃部分；
剩余部分在这样的集合上均为零，再由可数并处理 \(\sigma\)-有限集。
这一几何识别没有在本节证明，准确结论见前引 Remark 1.2 的式(1.1)--(1.4)。
结合\cref{b1:thm:bv-jump-structure}，
没有有限近似极限的集合除去 \(J_u\) 只有 \(\mathcal H^{d-1}\) 零测余项，
因此它也是 \(|D^cu|\) 的零测集。这给出最后一句。
\end{proofsketch}

\begin{proposition}\label{b1:prop:bv-decomposition-variation}
三部分分解唯一，并且对每个 Borel 集 \(A\subseteq\Omega\)，
\begin{equation}\label{b1:eq:bv-three-part-variation}
 |Du|(A)=\int_A|\nabla u|\,\mathrm dx
   +\int_{A\cap J_u}(u^+-u^-)\,\mathrm d\mathcal H^{d-1}
   +|D^cu|(A).
\end{equation}
\end{proposition}
\begin{proof}
Lebesgue 分解唯一确定第一项。
在 \(J_u\) 上限制剩余测度，唯一确定第二项；相减便唯一确定第三项。
取上述三个互不交的承载集，在每个承载集上 \(Du\) 等于相应分量，
所以其变差也等于该分量的变差。
由变差的可数可加性相加，再用向量密度的变差公式，即得所述等式。
\end{proof}

Cantor 部分的名称不是说它必须集中在三分 Cantor 集上，
而是说它代表不能由体积密度和跳跃面描述的奇异变化。
在一维，它恰好成为无原子的奇异部分，这一点可以不借高维结构定理直接证明。

\begin{proposition}\label{b1:prop:one-dimensional-bv-decomposition}
设 \(I=(a,b)\) 为有界区间，\(u\in BV(I)\)，\(\mu=Du\)。
则存在常数 \(c\)，使
\[
 \widehat u(x)=c+\mu\bigl((a,x]\bigr)
\]
是 \(u\) 的右连续代表。
其左右极限之差等于 \(\mu(\{x\})\)，且
\[
 D^ju=\sum_{x\in I}\mu(\{x\})\delta_x.
\]
这里非零项至多可数，绝对值之和有限。
测度 \(D^cu\) 则是 \(\mu\) 的奇异部分删去全部原子后的剩余测度；
它的累积分布函数连续。
\end{proposition}
\begin{proof}
函数 \(F(x)=\mu\bigl((a,x]\bigr)\) 有界。由符号测度积分的 Fubini 计算，
对 \(\varphi\in C_c^\infty(I)\)，
\[
 \int_a^b F(x)\varphi'(x)\,\mathrm dx
 =\int_I\left(\int_t^b\varphi'(x)\,\mathrm dx\right)\mathrm d\mu(t)
 =-\int_I\varphi\,\mathrm d\mu.
\]
因此 \(D(u-F)=0\)。
\secref{b1:sec:2004}证明的零弱导数判别说明 \(u-F=c\) 几乎处处。
有限测度的从上、从下连续性分别给出 \(F\) 的右连续性和左极限，
并且 \(F(x)-F(x-)=\mu(\{x\})\)。
非零原子至多可数，因为对每个正整数 \(n\)，质量绝对值至少 \(1/n\) 的原子只有有限个。
对有限原子集求和再取上确界，得到
\(\sum_x|\mu(\{x\})|\le|\mu|(I)\)。

累积分布代表在非原子点连续，在原子点具有不相同的两侧极限；
因此它的近似跳跃集恰为非零原子集。
每个原子的有向跳跃等于其有符号质量，给出 \(D^ju\) 的公式。
绝对连续测度没有原子，故原子全部属于 \(\mu\) 的奇异部分。
删去它们以后所得奇异测度无原子，其累积分布在每点左右连续，因而连续。
\end{proof}

例如，习题\ref{b1:ex:14-cantor-function}中的 \(F_C\) 满足
\[
 DF_C=\rho_C,\quad D^aF_C=0,\quad D^jF_C=0,\quad D^cF_C=\rho_C
 \quad\text{在 }(0,1)\text{ 上}.
\]
第一式由刚才的测试积分得出；后面三式来自 \(\rho_C\) 奇异且无原子。
于是函数
\[
 u(x)=x+\mathbf1_{(1/2,1)}(x)+F_C(x)
\]
在同一区间内同时具有三种非零变化：
\(D^au=m_1\)、\(D^ju=\delta_{1/2}\)、\(D^cu=\rho_C\)，总变差为 \(3\)。
这也表明连续性只能排除跳跃，不能排除 Cantor 部分。

最后记录复合函数所需的精细链律。
下一节会用它分别测试跳跃的存在与梯度的极限。

\begin{theorem}[Volpert 链律的 \(C^1\) 版本]
\label{b1:thm:volpert-chain-rule}
设 \(u\in BV(\Omega)\)，\(\Phi\in C^1(\mathbb R)\)，且 \(\Phi'\) 有界。
若 \(m_d(\Omega)<\infty\) 或 \(\Phi(0)=0\)，则 \(\Phi(u)\in BV(\Omega)\)，并且
\begin{equation}\label{b1:eq:volpert-chain-rule}
 D\bigl(\Phi(u)\bigr)
 =\Phi'(\widetilde u)(D^au+D^cu)
  +\bigl(\Phi(u^+)-\Phi(u^-)\bigr)\nu_u\,
    \mathcal H^{d-1}\mathbin{\llcorner}J_u.
\end{equation}
其中 \(\widetilde u\) 在有限近似极限不存在处可任意补值。
\end{theorem}
\term[volpert-lianlv]{Volpert 链律}[Volpert chain rule]在跳跃处使用两端值之差，
而不是把某个任选的点值代入 \(\Phi'\) 后乘上跳跃高度。
\begin{proofsketch}
一维时先用测度累积代表，把连续变化与原子变化分开：
连续部分满足通常的链式乘法，原子在复合后变成
\(\Phi(u(x+))-\Phi(u(x-))\)。
高维论证用有界变差函数的切片定理将各方向的分布导数还原为这些一维导数，
再重组为向量测度。这里省略的是切片定理及其与三个导数部分的相容性，
不是对公式作形式求导。
所用版本见 Alberti 与 Mantegazza 的《A Note on the Theory of SBV Functions》%
\cite[Remark 2.1, equation (2.1)]{AlbertiMantegazza1997SBV}；
本书按既定方向约定显式写出法向因子。
可积性由 \(|\Phi(t)|\le|\Phi(0)|+\|\Phi'\|_\infty\cdot|t|\) 及给定假设保证，
测度右侧的有限性则由三部分变差公式和 \(\Phi\) 的 Lipschitz 界保证。
\end{proofsketch}

% !TeX root = ../../Book1.tex
\section{特殊有界变差函数与紧性}
\label{b1:sec:2804}

在允许界面出现的变分问题中，变化的代价通常分成两类：
区域内部的梯度能量，以及界面本身的面积代价。
如果把所有有界变差函数都作为候选，而能量只计入这两项，
Cantor 部分便可能不付出任何代价。
因此需要一个排除该部分、同时保留跳跃的函数类。

\ProseParagraphBreak

本节先定义特殊有界变差函数，说明它对有界变差空间中的严格收敛并不闭合，
再记录带有值域、超线性梯度和跳跃面积控制的 Ambrosio 紧性定理。
定理用到\chapref{b1:ch:27}中有界变差函数的紧性，以及前节作为深层输入记录的 Volpert 链律。
后者如何用于特殊有界变差函数空间的闭合性，可参看 Alberti、Mantegazza 的
《A Note on the Theory of SBV Functions》%
\cite[Theorem 1.4, Propositions 2.3 and 2.5]{AlbertiMantegazza1997SBV}。
% 实读作者公开版第4--7页，对应刊印378--381页；
% Theorem1.4只陈述紧性及绝对连续/跳跃部分分开收敛，不包含面积下半连续性。
% 本节取其增长函数f(t)=1(t>0),f(0)=0，采用正弦/余弦测试展开闭合桥。

\subsection{特殊有界变差函数}
\label{b1:subsec:280401}

\begin{definition}\label{b1:def:sbv-function}
设 \(u\in BV(\Omega)\)。若 \(D^cu=0\)，则称 \(u\) 为%
\term[teshu-youjie-biancha-hanshu]{特殊有界变差函数}[special function of bounded variation]，
记 \(u\in SBV(\Omega)\)。\foreignidx{SBV}因此
\[
 Du=\nabla u\,m_d+(u^+-u^-)\nu_u\,
 \mathcal H^{d-1}\mathbin{\llcorner}J_u.
\]
\end{definition}
\symidx{\(SBV(\Omega)\)：特殊有界变差函数空间}

每个 \(W^{1,1}\) 函数都属于 \(SBV\)。
有界区域内有限周长集合的示性函数也属于 \(SBV\)：
它的导数由约化边界上的面积表示，故完全属于跳跃部分。
但特殊有界变差函数空间既不等于 Sobolev 空间，也不等于整个有界变差函数空间：
一个非零阶跃不属于 \(W^{1,1}\)，而前节的连续 Cantor 函数不属于 \(SBV\)。
定义本身也不要求跳跃集具有有限面积，更不要求它是闭集或由有限张光滑曲面组成。

\begin{proposition}\label{b1:prop:sbv-not-strictly-closed}
在区间 \((0,1)\) 上，特殊有界变差函数空间对有界变差空间中的严格收敛并不闭合。
甚至存在 \(0\le u_n\le1\) 的连续分段线性函数，
满足 \(u_n\to F_C\) 一致收敛、\(|Du_n|\bigl((0,1)\bigr)=1\)，
而 \(D^cF_C\ne0\)。
\end{proposition}
\begin{proof}
记 \(C_n\) 为三分 Cantor 集构造中的第 \(n\) 层闭区间之并。
在每个长度 \(3^{-n}\) 的基本区间上，用直线连接 Cantor 函数的两个端点值；
在其余间隙上保持 Cantor 函数的常值，所得函数记为 \(u_n\)。
两端值之差为 \(2^{-n}\)，故
\[
 u_n'=(3/2)^n\cdot\mathbf1_{C_n}\quad\text{几乎处处}.
\]
函数 \(u_n\) 连续且分段线性，因而属于 \(W^{1,1}\subseteq SBV\)。
它在每个基本区间内与 \(F_C\) 相差至多 \(2^{-n}\)，
在间隙上则相等，所以一致收敛成立。
又因为 \(m_1(C_n)=(2/3)^n\)，
\[
 \int_0^1|u_n'|\,\mathrm dx=1=|DF_C|\bigl((0,1)\bigr).
\]
这正是严格收敛，但 \(D^cF_C=\rho_C\ne0\)。
\end{proof}

这组函数没有跳跃，且梯度的 \(L^1\) 范数一致有界，
仍能把全部导数集中成奇异连续测度。
若 \(p>1\)，同一计算给出
\[
 \int_0^1|u_n'|^p\,\mathrm dx=(3/2)^{n(p-1)}\longrightarrow\infty.
\]
超线性的梯度控制正是后面紧性条件中的一个实质要求。

\subsection{紧性}
\label{b1:subsec:280403}

\begin{theorem}[Ambrosio 紧性：有界值域与超线性梯度（深层紧性定理）]
\label{b1:thm:ambrosio-sbv-compactness}
设 \(\Omega\) 为有界开集，\(1<p<\infty\)，\(u_j\in SBV(\Omega)\)，且
\[
 \sup_j\left\{\|u_j\|_{L^\infty(\Omega)}
       +\int_\Omega|\nabla u_j|^p\,\mathrm dx
       +\mathcal H^{d-1}(J_{u_j})\right\}<\infty.
\]
则存在子列与 \(u\in SBV(\Omega)\cap L^\infty(\Omega)\)，使
\[
 u_j\longrightarrow u\quad\text{在每个 }L^r(\Omega),\ 1\le r<\infty,
 \quad
 \nabla u_j\rightharpoonup\nabla u\quad\text{在 }L^p(\Omega;\mathbb R^d).
\]
此外，\(D^au_j\) 与 \(D^ju_j\) 分别\WeakStar{}收敛到 \(D^au\) 与 \(D^ju\)，
测度收敛均以 \(C_0(\Omega;\mathbb R^d)\) 为测试类。
\end{theorem}
\term[ambrosio-sbv-jinxing]{Ambrosio 紧性定理}[Ambrosio compactness theorem]
比一般有界变差函数的紧性多给出一个结论：极限仍没有 Cantor 部分。
\begin{proofsketch}
以下展开这个特殊版本的闭合步骤；所依赖的额外深层工具是
\cref{b1:thm:volpert-chain-rule}，不把它当作已经由普通 Sobolev 链律推出的结果。
记统一的值域界为 \(M\)，跳跃面积界为 \(S\)。
两侧近似值也在 \([-M,M]\) 中，故
\[
 |Du_j|(\Omega)
 \le m_d(\Omega)^{1-1/p}\|\nabla u_j\|_{L^p(\Omega)}
      +2M S.
\]
于是 \(BV\) 范数一致有界。由\cref{b1:thm:bv-local-compactness}取局部 \(L^1\) 收敛子列。
对于内部紧集穷竭 \(K_n\uparrow\Omega\)，
\(\int_{\Omega\setminus K_n}|u_j|\le M\,m_d(\Omega\setminus K_n)\) 一致趋零；
所以\cref{b1:cor:bv-global-compactness-with-tails}给出全域 \(L^1\) 收敛，
极限 \(u\in BV(\Omega)\) 且 \(|u|\le M\)。
再取子列，使 \(u_j\to u\) 几乎处处，且 \(\nabla u_j\rightharpoonup h\) 于 \(L^p\)。
尚需识别 \(h\) 并排除 \(D^cu\)。

取任意 \(\Phi\in C^1(\mathbb R)\)，满足 \(0\le\Phi\le1\) 且 \(\Phi'\) 有界。
由链律和 \(D^cu_j=0\)，测度
\[
 T_j^\Phi=D\bigl(\Phi(u_j)\bigr)-\Phi'(u_j)\nabla u_j\,m_d
\]
的总变差不超过 \(S\)。
函数 \(\Phi(u_j)\to\Phi(u)\) 于 \(L^1\)。
若 \(p'=p/(p-1)\)，对每个 \(q\in L^{p'}\)，
控制收敛给出
\(\Phi'(u_j)q\to\Phi'(u)q\) 于 \(L^{p'}\)，
因而 \(\Phi'(u_j)\nabla u_j\rightharpoonup\Phi'(u)h\)。
对紧支集光滑向量场取极限，再用变差的测试表述，得到
\[
 \left|D\bigl(\Phi(u)\bigr)-\Phi'(u)h\,m_d\right|(\Omega)\le S.
\]
再对 \(u\) 用链律，绝对连续、Cantor、跳跃三项互相奇异，因此
\[
 \int_\Omega|\Phi'(\widetilde u)|\,\mathrm d\sigma\le S,
 \quad
 \sigma=|\nabla u-h|\,m_d+|D^cu|.
\]
有限近似极限在 \(\sigma\)-几乎处处存在，故积分不依赖其余点的补值。
分别取
\[
 \Phi_N(t)=\frac{1+\sin(Nt)}2,\quad
 \Psi_N(t)=\frac{1+\cos(Nt)}2.
\]
因为 \(|\Phi_N'(t)|+|\Psi_N'(t)|\ge N/2\)，两次估计相加即得
\[
 \frac N2\,\sigma(\Omega)\le2S.
\]
令 \(N\to\infty\)，便有 \(\sigma=0\)，即 \(h=\nabla u\) 且 \(D^cu=0\)。
这个振荡测试正是所引论文 Propositions 2.3、2.5 在固定面积界下的核心机制。

最后，\(|u_j-u|\le2M\) 给出
\(\|u_j-u\|_r^r\le(2M)^{r-1}\|u_j-u\|_1\)，得到全部有限 \(r\) 的强收敛。
绝对连续部分的收敛来自梯度弱收敛，
跳跃部分的收敛来自 \(D^ju_j=Du_j-D^au_j\)
与\cref{b1:prop:bv-distributional-weak-star}。
\end{proofsketch}

这一定理为自由间断能量的直接法提供紧性，但尚未给出跳跃面积的下半连续性。
下一节先明确具体能量与容许类，再补上这项测度论接口。


\begin{chaptersummary}
BV 导数的三部分分别由体积、跳跃界面与不充余维一有限面积集的奇异变化承载。
这个区分不是把任意奇异测度机械拆成“有原子”和“无原子”：
只有在一维，原子恰与跳跃对应；高维的跳跃面测度本身也可以没有原子。

点值的恢复同样需要区分尺度。跳跃处没有单一有限近似极限，
但两个半球平均给出不同的有限值；它们的差与朝向共同决定跳跃导数。
给奇异测度乘一个函数时，应使用相应精细代表，
不能只凭 Lebesgue 几乎处处相等就任意替换。

SBV 排除了 Cantor 部分，却不自动对 BV 严格收敛闭合。
超线性梯度界与跳跃面积界分别防止不同形式的变化集中。
在有界值域内，这些界配合链律和下半连续性，给出 Mumford–Shah 型极小元。
存在性只产生函数等价类及其跳跃结构，没有同时给出唯一性、闭裂纹表示或界面正则性。
至此，本卷的积分、微分、覆盖、可整流性和紧性已经汇入同一个变分问题；
继续深入时，仍应逐项辨认哪些工具已经证明，哪些属于后续理论。
\end{chaptersummary}

\BookExercises

% 第二十八章习题临时稿；由根任务接在唯一的BookExercises之后，不另建入口。
% 八题均为自拟；所用BV结构、Volpert链律、SBV紧性来自本章明确陈述的版本，
% 不把其深层证明移入基础习题。逐题注释记录训练目标与相对正文的增量。

\begin{exercise}\label{b1:ex:28-radial-piecewise-decomposition}
在 \(\Omega=B(0,2)\subset\mathbb R^2\) 中，给定 \(a\in\mathbb R\)，令
\[
 u_a(x)=
 \begin{cases}
 |x|^2,&|x|<1,\\
 2|x|+a,&1<|x|<2.
 \end{cases}
\]
单位圆上的值任取。求 \(D^au_a,D^ju_a,D^cu_a\) 以及
\(|Du_a|(\Omega)\)，并判断 \(u_a\) 何时属于 \(W^{1,1}(\Omega)\)。
按照 \(u^+>u^-\) 的约定，分别写出可能跳跃处的两个值和法向。
\end{exercise}
% 来源：自拟；将两侧常数推广到两段径向光滑函数，训练体积分量、界面分量及排序反向。
\begin{solution}
两侧的经典梯度分别为 \(2x\) 和 \(2x/|x|\)。对内部光滑测试场逐区积分，
由\cref{b1:prop:smooth-domain-perimeter}，单位圆上的贡献等于外侧迹减内侧迹，
再乘从内向外的单位方向。因此
\[
 D^au_a=
 \left(2x\,\mathbf1_{\{|x|<1\}}
       +\frac{2x}{|x|}\mathbf1_{\{1<|x|<2\}}\right)m_2,
\]
\[
 D^ju_a=(a+1)x\,\mathcal H^1\llcorner\partial B(0,1),
 \quad D^cu_a=0.
\]
外圆是定义域的边界，不进入相对分布导数。两份非零分量互相奇异，所以
\[
 \begin{aligned}
 |Du_a|(\Omega)
 &=\int_{B(0,1)}2|x|\,\mathrm dx
   +\int_{B(0,2)\setminus B(0,1)}2\,\mathrm dx
   +2\pi|a+1|\\
 &=\frac{22\pi}{3}+2\pi|a+1|.
 \end{aligned}
\]
特别地，每个 \(u_a\) 都属于 \(SBV(\Omega)\)。

若 \(a>-1\)，圆周上的高值在外侧，
\[
 u_a^+=a+2,\quad u_a^-=1,\quad \nu_{u_a}=x.
\]
若 \(a<-1\)，高值改在内侧，
\[
 u_a^+=1,\quad u_a^-=a+2,\quad \nu_{u_a}=-x.
\]
两种情形中 \((u_a^+-u_a^-)\nu_{u_a}\) 都等于 \((a+1)x\)。
当 \(a=-1\) 时两侧迹相同，跳跃集为空；此时显式的有界梯度就是弱梯度，
甚至 \(u_{-1}\in W^{1,\infty}(\Omega)\)。
其余 \(a\) 有非零奇异导数，故不属于 \(W^{1,1}(\Omega)\)。
\end{solution}

\begin{exercise}\label{b1:ex:28-chain-jump-cancellation}
\begin{enumerate}
\item\label{b1:item:28-chain-square-cancellation}
在 \((-1,1)\) 上令 \(u(x)=x+\operatorname{sgn}x\)，原点任意补值。
取一个 \(C^1(\mathbb R)\) 函数 \(\Phi\)，使它在 \([-2,2]\) 上等于 \(t^2\)，且 \(\Phi'\) 有界。
计算 \(Du\) 与 \(D\bigl(\Phi(u)\bigr)\)，说明光滑的值域变换可以完全消除跳跃，却不使函数变成常数。
\item\label{b1:item:28-bilipschitz-value-change}
设 \(m_d(\Omega)<\infty\)、\(u\in BV(\Omega)\)，且
\[
 \Psi\in C^1(\mathbb R),\quad
 0<c\le|\Psi'(t)|\le C<\infty\quad(t\in\mathbb R).
\]
证明 \(J_{\Psi(u)}=J_u\)，写出排序后法向的变化，并证明
\[
 u\in SBV(\Omega)\quad\Longleftrightarrow\quad
 \Psi(u)\in SBV(\Omega).
\]
\end{enumerate}
\end{exercise}
% 来源：自拟；第一部分让两侧像值相合，第二部分辨明双Lipschitz值域变化不能产生同样消失。
\begin{solution}
第一部分中，原点两侧极限为 \(-1,1\)，所以
\[
 Du=m_1+2\delta_0.
\]
可具体选择 \(\Phi(t)=t^2\) 于 \(|t|\le2\)，在两侧分别以端点切线继续；
这样 \(\Phi\in C^1\) 且 \(|\Phi'|\le4\)。
复合函数为
\[
 \Phi(u(x))=(1+|x|)^2.
\]
它连续且 Lipschitz，其弱导数为
\[
 D\bigl(\Phi(u)\bigr)
 =2(1+|x|)\operatorname{sgn}x\,m_1.
\]
链律中的跳跃系数恰为 \(\Phi(1)-\Phi(-1)=0\)。
这不是把原点的代表值改掉，而是确实使两侧极限变为同一个值。

第二部分中，连续函数 \(\Psi'\) 从不为零，所以符号固定。
微积分基本定理给出
\[
 c|s-t|\le|\Psi(s)-\Psi(t)|\le C|s-t|.
\]
因此 \(\Psi\) 单调且映满 \(\mathbb R\)，其逆也为 Lipschitz 函数。
半球中的 \(L^1\) 逼近误差在复合后至多乘以 \(C\)，逆向至多乘以 \(1/c\)；
所以两侧不同有限值的跳跃点恰好保持不变。
若 \(\Psi'>0\)，高低次序与 \(\nu_u\) 不变；若 \(\Psi'<0\)，
高低次序交换，\(\nu_{\Psi(u)}=-\nu_u\)。
两种情形的跳跃高度均为 \(|\Psi(u^+)-\Psi(u^-)|\)。

由\cref{b1:thm:volpert-chain-rule}，
\[
 D^c\bigl(\Psi(u)\bigr)=\Psi'(\widetilde u)\,D^cu.
\]
近似代表在 \(|D^cu|\)-几乎处处有限，故
\[
 c|D^cu|\le |D^c\bigl(\Psi(u)\bigr)|\le C|D^cu|.
\]
Cantor 部分为零当且仅当复合后的 Cantor 部分为零，得到 \(SBV\) 的等价性。
有限体积条件保证即使 \(\Psi(0)\ne0\)，复合函数仍属于 \(L^1(\Omega)\)。
\end{solution}

\begin{exercise}\label{b1:ex:28-cantor-cylinder-and-atoms}
设 \(\mathcal C\subset[0,1]\) 为三分 Cantor 集，\(C\) 为其连续分布函数，
\(\mu_C=DC\) 为无原子的 Cantor 概率测度。在 \(\Omega=(0,1)^2\) 中令
\[
 u(x,y)=C(x),\quad v(x,y)=\mathbf1_{\{y>1/2\}}.
\]
\begin{enumerate}
\item\label{b1:item:28-cantor-cylinder-decomposition}
计算二者的三分解，并说明 \(|Du|\)、\(|Dv|\) 都没有原子，
但一份是纯 Cantor 部分，另一份是纯跳跃部分。
\item\label{b1:item:28-cantor-cylinder-diffuse}
不调用 Cantor 部分的结构定理，直接证明
\[
 (\mu_C\otimes m_1)(A)=0
\]
对每个 \(\mathcal H^1\)-\(\sigma\) 有限的 Borel 集 \(A\subset\Omega\) 成立。
据此说明 \((\mathcal C\times(0,1))\cap\Omega\) 不是
\(\mathcal H^1\)-\(\sigma\) 有限集。
\end{enumerate}
\end{exercise}
% 来源：自拟；从一维Cantor函数增至二维乘积，并用互不相交的竖截面直接区分无原子与无跳跃。
\begin{solution}
对测试函数先沿 \(x\) 积分，再沿 \(y\) 积分，得到
\[
 Du=e_1(\mu_C\otimes m_1),\quad
 Dv=e_2\mathcal H^1\llcorner\bigl((0,1)\times\{1/2\}\bigr).
\]
第一份测度集中在体积为零的 Cantor 柱上。函数 \(u\) 连续，没有跳跃点，
所以 \(D^cu=Du\)，而 \(D^au=D^ju=0\)。
函数 \(v\) 在水平线两侧有不同常值，因此 \(D^jv=Dv\)，其他两部分为零。
乘积测度和线上的长度测度都不给单点质量。
可见高维导数测度“无原子”并不等于“没有跳跃部分”。

先设 \(\mathcal H^1(A)<\infty\)。对每个 \(x\)，记
\(A_x=\{y:(x,y)\in A\}\)。
不同竖线上的集合互不相交，且由等距不变性，
\[
 \mathcal H^1\bigl(\{x\}\times A_x\bigr)=m_1(A_x).
\]
有限测度空间中，互不相交且各自具有正测度的集合至多可数：
对每个正整数 \(n\)，其中测度至少为 \(1/n\) 的集合只能有有限个。
所以 \(m_1(A_x)>0\) 的 \(x\) 至多可数。
由 Tonelli 定理与 \(\mu_C\) 无原子，
\[
 (\mu_C\otimes m_1)(A)=\int m_1(A_x)\,\mathrm d\mu_C(x)=0.
\]
对 \(\mathcal H^1\)-\(\sigma\) 有限集作可数分解，即得一般结论。
另一方面，\(\mu_C\otimes m_1\) 给 Cantor 柱的质量为 \(1\)，
所以该柱不可能具有所述 \(\sigma\) 有限性。
\end{solution}

\begin{exercise}\label{b1:ex:28-jump-cantor-approximation}
沿用习题\ref{b1:ex:28-cantor-cylinder-and-atoms}的一维 Cantor 分布函数。
令 \(I_{n,k}\)、\(1\le k\le2^n\)，为第 \(n\) 步留下的基本闭区间，
按从左到右排列，并记其中心为 \(c_{n,k}\)。定义
\[
 \mu_n=2^{-n}\sum_{k=1}^{2^n}\delta_{c_{n,k}},
 \quad
 u_n(x)=\mu_n((0,x])\quad(0<x<1).
\]
证明 \(u_n\in SBV(0,1)\)，且 \(u_n\to C\) 严格于 \(BV(0,1)\)，
但极限不属于 \(SBV(0,1)\)。
分别计算
\[
 \mathcal H^0(J_{u_n}),\quad
 \int_{J_{u_n}}(u_n^+-u_n^-)^q\,\mathrm d\mathcal H^0
 \quad(q>0),
\]
并指出\cref{b1:thm:ambrosio-sbv-compactness}中的哪项统一控制失效。
最后判断这是否也是通常 \(BV\) 范数下的收敛。
\end{exercise}
% 来源：自拟；与正文连续分段线性逼近相对照，令梯度恒为零而跳跃点数量发散，并量化振幅罚项。
\begin{solution}
每个 \(u_n\) 是有限阶跃函数，\(Du_n=\mu_n\)，所以它只有跳跃部分，
\[
 |Du_n|\bigl((0,1)\bigr)=1,\quad \nabla u_n=0
 \quad\text{几乎处处}.
\]
每个基本区间的 Cantor 质量与 \(\mu_n\) 质量都等于 \(2^{-n}\)。
在基本区间之间的间隙，两份分布函数取同一个值；在一个基本区间内，
二者都位于相同的两个相邻累计质量之间。因此
\[
 \sup_{0<x<1}|u_n-C|\le2^{-n}.
\]
于是有 \(L^1\) 收敛，同时 \(|Du_n|\bigl((0,1)\bigr)=|DC|\bigl((0,1)\bigr)=1\)，
这正是严格收敛。极限连续而具有非零奇异导数，故其导数全是 Cantor 部分，不属于 \(SBV\)。

每个中心点的跳跃高度为 \(2^{-n}\)，故
\[
 \mathcal H^0(J_{u_n})=2^n,\quad
 \int_{J_{u_n}}(u_n^+-u_n^-)^q\,\mathrm d\mathcal H^0
 =2^{n(1-q)}.
\]
这列函数的值域一致受控，所有 \(p>1\) 的梯度积分都为零；
失效的是跳跃集的面积，即此处的点数控制。
特别地，\(q\ge1\) 时振幅的 \(q\) 次罚项有界，仍不足以代替跳跃点数；
\(q>1\) 时它甚至趋零。

最后，有限原子测度 \(\mu_n\) 与无原子的 \(\mu_C\) 相互奇异，所以
\[
 |D(u_n-C)|\bigl((0,1)\bigr)=|\mu_n-\mu_C|\bigl((0,1)\bigr)=2.
\]
因此并没有通常 \(BV\) 范数的强收敛。严格收敛只要求各自总变差的数值收敛，
不要求导数测度之差的全变差趋零。
\end{solution}

\begin{exercise}\label{b1:ex:28-bv-removable-interface}
设 \(d\ge2\)，\(K\) 是开集 \(\Omega\) 中的相对闭集，并且
\(\mathcal H^{d-1}|_K\) 为 \(\sigma\) 有限。
给定 \(u\in BV(\Omega)\)，假设
\[
 u\in W_{\mathrm{loc}}^{1,1}(\Omega\setminus K),
 \quad \mathcal H^{d-1}(K\cap J_u)=0.
\]
证明 \(u\in W^{1,1}(\Omega)\)。
分别利用习题\ref{b1:ex:28-radial-piecewise-decomposition}和%
\ref{b1:ex:28-cantor-cylinder-and-atoms}说明：
若删除“没有界面跳跃”的条件，或删除 \(K\) 的 \(\sigma\) 有限条件，结论均可能失败。
\end{exercise}
% 来源：自拟；将三分解转为带BV先验的奇集可去性，两个反例分别孤立跳跃和Cantor障碍。
\begin{solution}
在开集 \(\Omega\setminus K\) 内，局部 Sobolev 条件说明分布导数由局部可积函数表示。
按可数内部开覆盖及分解唯一性，\(D^ju\)、\(D^cu\) 在该开集上的限制均为零，
所以两份奇异测度都集中在 \(K\) 上。

跳跃表示与 \(\mathcal H^{d-1}(K\cap J_u)=0\) 给出 \(|D^ju|(K)=0\)。
由\cref{b1:thm:bv-derivative-decomposition}中 Cantor 部分不充
\(\mathcal H^{d-1}\)-\(\sigma\) 有限集的结论，\(|D^cu|(K)=0\)。
因此 \(D^su=0\)。再由\cref{b1:prop:bv-sobolev-singular-criterion}，
得到 \(u\in W^{1,1}(\Omega)\)。
这里的全域可积性来自原先的 \(BV(\Omega)\) 假设，不是仅由域外局部 Sobolev 性推出。

若删除跳跃条件，取第一题中 \(a\ne-1\) 的函数，并令 \(K=\partial B(0,1)\)。
它在 \(K\) 外分片光滑，\(K\) 的长度有限，却有非零界面导数。
若删除 \(\sigma\) 有限条件，取第三题中的 \(u(x,y)=C(x)\) 及 Cantor 柱 \(K\)。
它连续，故没有跳跃；在 \(K\) 外局部常于 \(x\)，因而光滑，
却仍有非零 Cantor 导数。第三题已经证明这时 \(K\) 确实不是长度的 \(\sigma\) 有限集。
\end{solution}

\begin{exercise}\label{b1:ex:28-cantor-reparameterization}
设 \(C\) 与 \(\mu_C\) 如前，\(\Phi\in C^1(\mathbb R)\) 且 \(\Phi'\) 有界。
证明
\[
 C_\#\mu_C=m_1\llcorner[0,1],
 \quad
 |D^c\bigl(\Phi(C)\bigr)|((0,1))=\int_0^1|\Phi'(t)|\,\mathrm dt.
\]
据此证明：\(\Phi(C)\in SBV(0,1)\) 当且仅当 \(\Phi\) 在 \([0,1]\) 上为常数。
将这一结论与习题\ref{b1:ex:28-chain-jump-cancellation}中的消跳现象比较。
\end{exercise}
% 来源：自拟；通过分布函数的推前测度精确计算链律中的Cantor能量，不重复原始Cantor反例。
\begin{solution}
对 \(0<t<1\)，由 \(C\) 连续、非减且映满 \([0,1]\)，
存在最大的 \(b_t\in[0,1]\)，使 \(C(b_t)=t\)。
于是
\[
 \{x:C(x)\le t\}=[0,b_t],\quad
 \mu_C(\{C\le t\})=\mu_C([0,b_t])=C(b_t)=t.
\]
端点 \(t=0,1\) 同样成立。比较区间的分布函数并用测度唯一性，
得到 \(C_\#\mu_C=m_1\llcorner[0,1]\)。

函数 \(C\) 连续，近似代表就是它本身，导数为纯 Cantor 测度。
由链律，
\[
 D\bigl(\Phi(C)\bigr)=\Phi'(C)\,\mu_C.
\]
复合函数仍连续，因而没有跳跃；上式又集中在 Lebesgue 零集上，
所以整份导数都是 Cantor 部分。推前等式给出
\[
 |D^c\bigl(\Phi(C)\bigr)|((0,1))
 =\int|\Phi'(C(x))|\,\mathrm d\mu_C(x)
 =\int_0^1|\Phi'(t)|\,\mathrm dt.
\]
它为零当且仅当连续函数 \(\Phi'\) 在 \([0,1]\) 上恒为零，
也就是 \(\Phi\) 在该区间为常数。
跳跃只比较两个端点的像值，两点可以在非恒定变换下合并；
这里的奇异变化却经 \(C\) 遍历整个值域，不能只让两个值相合就消除。
\end{solution}

\begin{exercise}\label{b1:ex:28-energy-scaling}
设 \(\Omega\subset\mathbb R^d\) 有界，\(u\in BV(\Omega)\)。
给定 \(\rho>0\)、\(s\ne0\)、\(b\in\mathbb R\)，令
\[
 \Omega_\rho=\rho\Omega,\quad
 T_\rho(x)=\rho x,\quad
 v(x)=s\,u(x/\rho)+b.
\]
证明三份导数均满足
\[
 D^\ell v=s\rho^{d-1}(T_\rho)_\#D^\ell u,
 \quad \ell\in\{a,j,c\},
\]
并写出跳跃集与法向的变化。
进一步设 \(u\in SBV(\Omega)\)、\(g\in L^2(\Omega)\)、\(1<p<\infty\)，
给定 \(\beta,\lambda>0\)，记允许取值 \(+\infty\) 的能量为
\[
 \mathcal F_{\beta,\lambda}(u;\Omega,g)
 =\int_\Omega|\nabla u|^p\,\mathrm dx
 +\beta\mathcal H^{d-1}(J_u)
 +\lambda\int_\Omega|u-g|^2\,\mathrm dx.
\]
令 \(g_\rho(x)=s\,g(x/\rho)+b\)，求这三项各自的尺度。
再求 \(\beta_\rho,\lambda_\rho\)，使
\[
 \mathcal F_{\beta_\rho,\lambda_\rho}(v;\Omega_\rho,g_\rho)
 =|s|^p\rho^{d-p}\mathcal F_{\beta,\lambda}(u;\Omega,g).
\]
\end{exercise}
% 来源：自拟；同时追踪体积、余维一面积及振幅，要求权重随尺度改变而不是只比较导数范数。
\begin{solution}
在分布测试中作 \(x=\rho y\) 的变量代换，得到
\[
 Dv=s\rho^{d-1}(T_\rho)_\#Du.
\]
伸缩保持 Lebesgue 零集和正则分布密度的类型，也把半球近似跳跃条件逐项带到新点，
因此
\[
 J_v=\rho J_u,\quad
 \nu_v(\rho x)=\operatorname{sgn}(s)\nu_u(x).
\]
跳跃值是 \(su^\pm+b\) 重新排序后的两个值。绝对连续部分和跳跃部分分别变换为题述形式，
从总导数中减去这两部分，Cantor 部分也得到同一公式。
特别地，
\[
 \nabla v(x)=\frac{s}{\rho}\nabla u(x/\rho)
 \quad\text{几乎处处}.
\]

于是三项分别为
\[
 \int_{\Omega_\rho}|\nabla v|^p
 =|s|^p\rho^{d-p}\int_\Omega|\nabla u|^p,
\]
\[
 \mathcal H^{d-1}(J_v)=\rho^{d-1}\mathcal H^{d-1}(J_u),
 \quad
 \int_{\Omega_\rho}|v-g_\rho|^2
 =s^2\rho^d\int_\Omega|u-g|^2.
\]
比较三个系数，得到
\[
 \beta_\rho=\beta|s|^p\rho^{1-p},
 \quad
 \lambda_\rho=\lambda|s|^{p-2}\rho^{-p}.
\]
可见梯度、跳跃面积和拟合误差不具有相同的长度尺度。
若同时保留正文的值域约束，原来的 \(|u|\le M\) 应相应改为
\(|v|\le |s|M+|b|\)；不能在任意振幅变换后仍声称同一个固定值域类不变。
\end{solution}

\begin{exercise}\label{b1:ex:28-one-jump-mumford-shah}
取 \(L,a,\beta,\lambda>0\)，在 \(\Omega=(-L,L)\) 上令
\(g(x)=a\operatorname{sgn}x\)，并考虑 \(p=2\) 的能量
\[
 \mathcal F(u)=\int_{-L}^L|u'|^2\,\mathrm dx
 +\beta\mathcal H^0(J_u)
 +\lambda\int_{-L}^L|u-g|^2\,\mathrm dx.
\]
先把容许函数限制为 \([-a,a]\) 值的常数函数，或至多有一个跳跃的分段常值函数。
求此受限类中的最小值及全部极小元，说明何时出现两个不同的极小元。
然后用一个无跳跃的线性函数说明：这个阈值不能直接宣称为完整
\(SBV\) 容许类中的全局阈值。
\end{exercise}
% 来源：自拟；精确解一个有限参数模型，再用新竞争者主动检验“受限极小值”与真正全局极小值的差别。
\begin{solution}
常数函数的最优值为 \(0\)，因为
\[
 \lambda\int_{-L}^L|c-g|^2\,\mathrm dx
 =2\lambda L(c^2+a^2).
\]
这给出能量 \(2\lambda La^2\)。

若实际跳跃发生在 \(t\in[0,L)\)，左右两段分别以数据的均值作为最优常数：
\[
 c_1=a\frac{t-L}{t+L},\quad c_2=a.
\]
二者都在容许值域内且彼此不同。右段没有拟合误差，左段的最小平方误差为
\[
 a^2(t+L)-a^2\frac{(t-L)^2}{t+L}
 =\frac{4a^2Lt}{L+t}.
\]
若 \(t<0\)，相同计算给出 \(4a^2L|t|/(L+|t|)\)。
因此一个真实跳跃的最小能量是
\[
 \beta+\frac{4\lambda a^2L|t|}{L+|t|}\ge\beta,
\]
且等号只在 \(t=0\) 时成立，此时函数正是 \(g\)。
故受限类的最小值为
\[
 \min\{\beta,\,2\lambda La^2\}.
\]
若 \(\beta<2\lambda La^2\)，唯一极小元为 \(g\)；
若 \(\beta>2\lambda La^2\)，唯一极小元为零函数；
若二者相等，这两个函数恰为全部极小元。函数在单点上的重定义不计作新的极小元。

现在取 \(u_\varepsilon(x)=\varepsilon x\)，其中 \(0<\varepsilon L\le a\)。
它满足同一值域约束且没有跳跃，但不属于上面的分段常值受限类。直接积分得到
\[
 \mathcal F(u_\varepsilon)-\mathcal F(0)
 =-2\lambda aL^2\varepsilon
  +\left(2L+\frac{2\lambda L^3}{3}\right)\varepsilon^2.
\]
充分小的正 \(\varepsilon\) 使右边为负。
特别在 \(\beta=2\lambda La^2\) 时，它的能量同时低于上述两个受限极小元。
所以受限类中的非唯一性计算完全成立，却没有证明完整模型在该参数处具有同样的极小元或阈值。
\end{solution}
